% ─── Chapter 4: Reinforcement Learning Framework ──────────────────
\chapter{Reinforcement Learning Framework}
\label{ch:rl_framework}

\section{MDP Formulation}

The Pok\'{e}mon battle environment is modelled as a Markov Decision Process
(MDP) with the following components:

\begin{center}
\small
\begin{tabularx}{\textwidth}{l X}
  \toprule
  \textbf{Component} & \textbf{Definition} \\
  \midrule
  State ($s$)       & Dictionary observation: token matrix $(13 \times 164)$,
                      species/item/ability ID vectors $(13,)$ each, action
                      mask $(14,)$.  Full specification in
                      Chapter~\ref{ch:architecture}. \\
  Action ($a$)       & Discrete choice from 14 possible actions (4 moves,
                      4 gimmick-enhanced moves, 6 switches), filtered by the
                      action mask to only legal options. \\
  Reward ($r$)       & Shaped reward computed each step from HP changes,
                      fainting events, and a terminal outcome bonus/penalty.
                      Detailed in Section~\ref{sec:reward_design}. \\
  Transition         & Deterministic given both players' actions;
                      stochasticity arises from the opponent's policy and
                      damage variance in the simulator. \\
  Discount ($\gamma$) & 0.97 based on a hyper parameter sweep. \\
  \bottomrule
\end{tabularx}
\end{center}

An episode corresponds to a single battle.  It terminates when one side has no
remaining conscious Pok\'{e}mon (win or loss).  There is no explicit time limit
or truncation.

\section{PPO Algorithm}

The agent is trained with \textbf{Proximal Policy Optimization} (PPO), an
on-policy actor-critic algorithm that balances sample efficiency with training
stability through a clipped surrogate objective.  PPO was chosen for the
following reasons:

\begin{itemize}
  \item \textbf{Stability:} The clipped objective prevents destructively large
        policy updates, which is critical when learning in a sparse-reward,
        high-variance environment.
  \item \textbf{Distributed rollouts:} PPO naturally integrates with Ray
        RLlib's distributed rollout architecture, enabling many parallel
        battles.
  \item \textbf{Action masking:} PPO's categorical distribution can be
        combined with action masking (setting invalid-action logits to a large
        negative value) without modifying the algorithm itself.
\end{itemize}

The policy network (actor) outputs a categorical distribution over the 14
actions, and the value network (critic) outputs a scalar state-value estimate
$V(s)$.  Both share a common transformer trunk (see
Chapter~\ref{ch:architecture}).

\section{Reward Design}
\label{sec:reward_design}

The reward function combines several components to provide dense learning signal
while preserving the primacy of winning.  All rewards are multiplied by a
global \texttt{reward\_scale}$\,=0.1$ before being returned to the agent, which
scales returns from approximately $[-15, +15]$ to $[-1.5, +1.5]$.  This scaling
was critical for value-function learning (see
Section~\ref{sec:reward_scaling}).

\begin{center}
\small
\begin{tabularx}{\textwidth}{l r X}
  \toprule
  \textbf{Component} & \textbf{Default} & \textbf{Description} \\
  \midrule
  Victory reward     & $+10.0$  & Agent wins the battle (terminal step). \\
  Defeat penalty     & $-10.0$  & Agent loses the battle (terminal step). \\
  HP differential    & $\times 2.0$ &
    $(\text{our\_team\_HP} - \text{opp\_team\_HP}) \times w$ as a shaped
    step reward.  Positive when the agent is ahead, negative when behind. \\
  Opponent fainted   & $+5.0$   & Per opponent Pok\'{e}mon knocked out. \\
  Own fainted        & $-5.0$   & Per agent Pok\'{e}mon that faints. \\
  Step penalty       & $-0.005$ & Per step, encourages battle efficiency
                                    (fewer turns). \\
  Matchup quality    & $\times 0.2$ &
    $+1.0$ if a super-effective move is available, $-0.5$ if only resisted
    moves, $-1.0$ if immune.  Rewards good type-readiness. \\
  Action quality     & $\times 0.3$ &
    Penalty for clearly suboptimal moves; bonus for resisting the opponent's
    best move. \\
  \bottomrule
\end{tabularx}
\end{center}

The reward is computed each step by \texttt{poke-env}'s
\texttt{reward\_computing\_helper}, parameterised by the active
\texttt{RewardConfig}.  At the terminal step, explicit outcome rewards are
applied on top.  Reward parameters are adjusted across curriculum stages (see
below). These were gradually introduced during development in descending order.

\section{Curriculum Learning}
\label{sec:curriculum}

Due to the complexity of the full challenge, training follows a \textbf{curriculum
learning} approach where the agent is first exposed to easy opponents and
gradually faces stronger ones.  The curriculum currently has four stages:

\begin{enumerate}
  \item \textbf{Moves \& Switches Warmup.}
        Opponent mix: 40\,\% random, 30\,\% \texttt{random\_no\_switch},
        30\,\% self-play.  Promotion threshold: 65\,\% win rate.
        Asymmetric rewards (victory $+8$, defeat $-10$) encourage caution.
        Higher HP weight (3.0) and faint values ($\pm 6.0$) provide strong
        per-step signal for learning basic move selection.

  \item \textbf{Random, More Moves.}
        Opponent mix: 70\,\% self-play, 10\,\% random,
        20\,\% \texttt{random\_no\_switch}.  Promotion threshold: 20\,\% win
        rate (intentionally low: the goal is exposure to self-play, not
        mastery).  Victory reward raised to $+10$.  Large minimum sample
        requirement (5\,000 episodes) ensures sufficient exposure before
        progressing.

  \item \textbf{Self-Play.}
        Opponent mix: 60\,\% \texttt{random\_no\_switch},
        40\,\% random.  Promotion threshold: 70\,\% win rate.
        After the self-play exposure stage, the agent trains against
        non-switching random opponents to solidify move selection skills.
        Action quality weight increased to 0.4 to refine move choice.

  \item \textbf{Mixed Final.}
        Opponent mix: 50\,\% heuristic, 30\,\% self-play,
        20\,\% \texttt{random\_no\_switch}.  Terminal stage (promotion
        threshold 101\,\%, i.e.\ never promotes).  The agent faces the
        strongest opponent mix, including the heuristic AI which considers
        type effectiveness and switches strategically.
\end{enumerate}

% TODO: If a fifth stage is added, update the list above.

\subsection{Per-Stage Reward Variants}

Reward parameters are adjusted across stages to reflect changing difficulty:

\begin{center}
\small
\begin{tabularx}{\textwidth}{l c c c c}
  \toprule
  \textbf{Parameter} & \textbf{Stage 1} & \textbf{Stage 2} & \textbf{Stage 3} & \textbf{Stage 4} \\
  \midrule
  Victory / Defeat   & $+8$ / $-10$ & $\pm 10$ & $\pm 10$ & $\pm 10$ \\
  HP weight          & 3.0   & 3.0   & 3.0   & 3.0   \\
  Opp.\ fainted      & $+6.0$ & $+6.0$ & $+6.0$ & $+6.0$ \\
  Own fainted        & $-6.0$ & $-6.0$ & $-6.0$ & $-6.0$ \\
  Step penalty       & $-0.01$ & $-0.01$ & $-0.01$ & $-0.01$ \\
  Matchup weight     & 0.1   & 0.1   & 0.1   & 0.1   \\
  Action quality wt. & 0.3   & 0.3   & 0.4   & 0.4   \\
  \bottomrule
\end{tabularx}
\end{center}

Note that Stage~1 uses asymmetric rewards (victory $+8$ vs.\ defeat $-10$) to
discourage early aggression, and a higher step penalty ($-0.01$) to encourage
efficient play from the start.

Promotion decisions use a rolling window of the last 300~episodes (minimum 50
samples required before promotion is considered; minimum 300 episodes must
elapse before any promotion). Total possible rewards is by design made in a way that can only increase with difficulty.

\section{Self-Play Mechanism}
\label{sec:self_play}

Self-play is implemented through a \texttt{SelfPlayPlayer} class that:

\begin{enumerate}
  \item Loads the latest model weights from
        \texttt{checkpoints/selfplay\_latest.pt} on first use.
  \item Re-checks the weights file each turn, enabling automatic updates
        when new checkpoints are saved.
  \item Maintains per-battle LSTM state for temporal memory.
  \item Converts the agent's compressed 14-action output to native
        \texttt{poke-env} \texttt{BattleOrder} objects.
  \item Falls back to \texttt{RandomPlayer} behaviour on inference errors.
\end{enumerate}

At each checkpoint, the training pipeline exports the current model weights to
the self-play weights path, making them immediately available to all
\texttt{SelfPlayPlayer} instances across workers.  This creates a natural
curriculum where the agent continuously adapts to an opponent of increasing
skill.

% TODO: Add talking point about self-play training results once available.
% Specifically: win rate curves for self-play stage, comparison with
% non-self-play training, and any observed oscillation or instability.

\section{Opponent Types}

The environment supports four opponent categories:

\begin{center}
\small
\begin{tabularx}{\textwidth}{l l X}
  \toprule
  \textbf{Key} & \textbf{Player Class} & \textbf{Behaviour} \\
  \midrule
  \texttt{random\_no\_switch} & \texttt{RandomNoSwitchPlayer} &
    Uniform random among legal \emph{move} orders; switches only when forced
    (no moves available or forced by simulator). \\
  \texttt{random} & \texttt{RandomPlayer} &
    Uniform random over all legal orders (moves \emph{and} switches). \\
  \texttt{heuristic} & \texttt{SimpleHeuristicsPlayer} &
    Rule-based AI: considers type effectiveness, switches when at a
    disadvantage. \\
  \texttt{self} & \texttt{SelfPlayPlayer} &
    Uses the training model's weights for inference with hot-reload support. \\
  \bottomrule
\end{tabularx}
\end{center}

In the embedding layer, \texttt{random\_no\_switch} is mapped to the same
bucket as \texttt{random} for the opponent-type feature, avoiding distribution
shift at the observation level when mixing or swapping baselines. From best to worst in terms of performance: heuristic, random no switch and random.
